\chapter*{Abstract}

Graphs arise in many scientific and computational domains, ranging from social networks to biology and chemistry. Their
irregular, non-Euclidean structure makes them expressive, but also difficult to
represent, generate, and evaluate. This thesis develops robust and
interpretable methods for graph representation, conditional graph generation,
and statistically grounded evaluation of graph generative models.

The first contribution is the Neighborhood Subgraph Pairwise Path Kernel
(NSPPK), an explicit graph kernel for attributed graphs. NSPPK represents a
graph using pairs of local subgraphs centred on selected anchor nodes, together
with the graph regions that connect those anchors. The connector subgraphs are
constructed from unions of shortest paths between the anchor nodes, allowing
the kernel to capture both the structure surrounding each anchor and the
structure linking different parts of the graph. NSPPK also incorporates
continuous node attributes directly into its feature map, avoiding
discretisation while preserving a sparse, deterministic, and interpretable
representation. Experiments show that NSPPK performs competitively with
established graph kernels and graph neural networks while offering strong
sample efficiency and computational transparency.

The second contribution is the Graph-Conditioned Diffusion Generator (GCDG), a
conditional diffusion-based graph generator that uses graph-level feature
vectors to guide the generation of node-level features before constructing a
final discrete graph. Its central focus is the graph-level-to-node-level
generation problem: translating a global graph descriptor into node existence,
node labels, degree information, and pairwise edge structure. The model
separates continuous diffusion-based denoising, node and edge prediction, and
constraint-aware decoding. This design makes intermediate predictions
inspectable and allows graph feasibility constraints to be enforced explicitly
during decoding. Experiments on synthetic graph families and QM9-like molecular
graphs show that conditional graph generation should be evaluated not only by
the realism of the generated graphs, but also by how faithfully they preserve
the information encoded in the conditioning vectors.

The third contribution is RankGen, a statistically grounded framework for
generative model evaluation. RankGen treats evaluation as a set of predictive
probes rather than as a single scalar similarity score. It defines
complementary metrics for Quality, Utility, Indistinguishability, and
Similarity, combines them with PAC-style reliability bounds, and ranks
generators using robust quantile-based summaries. Experiments across
controlled perturbations, molecular graph generators, image benchmarks, and
open-domain text generation demonstrate that RankGen can diagnose failure
modes such as memorisation, mode collapse, poor augmentation value, and
distributional mismatch.

Together, these contributions support a unified view of graph generative
modelling: reliable generation requires graph representations that expose
meaningful structure, generative models whose outputs can be explicitly
conditioned, and evaluation protocols that distinguish genuine task-relevant
utility from superficial structural similarity.
